% AUTO-GENERATED by tests/perf/compare_perf.py
\begin{table}[t]
\centering
\begingroup
\footnotesize
\setlength{\tabcolsep}{2pt}
\begin{tabular}{@{}p{0.10\linewidth}p{0.20\linewidth}p{0.13\linewidth}r r r p{0.14\linewidth}@{}}
\toprule
Module & Estimator & Reference & Max. $N$ & sp & Ref. & Faster path \\
\midrule
\code{01\_hdfe} & HDFE 2-way FE & \pkg{fixest} & 1{,}000{,}000 & 0.188 & 0.088 & R $\mathbf{2.1\times}$ \\
\code{02\_csdid} & CS-DiD simple ATT & \pkg{did} & 125{,}000 & 0.033 & 0.241 & sp $\mathbf{7.4\times}$ \\
\code{03\_scm} & Classical SCM & \pkg{Synth} & 100 & 11.636 & 2.411 & R $\mathbf{4.8\times}$ \\
\code{04\_dml} & DML PLR (lin. learners) & \pkg{DoubleML} Python & 10{,}000 & 0.008 & 0.009 & tie $1.12\times$ \\
\bottomrule
\end{tabular}
\endgroup
\caption{Track C measured performance on Apple Silicon arm64 / 8 cores / 24~GB. Times are median seconds at the largest sample size, across 3 or 5 repetitions after one warmup, using seed-fixed DGPs. Bold denotes a gap of at least $1.5\times$. The DML row uses \pkg{DoubleML} Python with the same linear learners, five folds, and public-workflow construction cost. The Callaway--Sant'Anna row uses a homogeneous-effect DGP favourable to the vectorised group-time loop.}
\label{tab:track-c-perf}
\end{table}
